\section{Measurement critique: TF-IDF as a representation, not an instrument}
\label{sec:measurement}

TF-IDF was designed for information retrieval \citep{sparckjones1972,manning2008}.
Retrieval ranks documents against a query, so a document's absolute scale is a nuisance
parameter: what matters is whether this document is a better match than that one, and a
report that says the same things twice as often is not a better match. Measurement asks a
different question --- how much of a construct a document contains --- and there the scale
\emph{is} the quantity of interest. The object of this section is the representation as it
comes out of the box: idf term weights, and the unit-length scaling of the document vector
that is the library default. Both are calibrated for the first question and neither is safe
for the second, but they are not equally unsafe. \textbf{The vector normalisation is
first-order and the term weighting is second-order}, and the section is ordered accordingly.
The argument concerns a class of methods, it holds analytically, and it would hold unchanged
if the original had documented every setting it used. It is not a claim about the original's
conclusion; that claim belongs to \S\ref{sec:artefact} and rests on evidence of a different
kind.

One attribution point governs everything below, and we concede it at the outset rather than
leave a referee to raise it. The original \citep{lagasio2024} discloses four
\texttt{TfidfVectorizer} settings and never mentions \texttt{norm} (\S\ref{sec:original}).
The four it discloses are exactly four non-default arguments, so a reader following the
ordinary convention that only departures from library defaults are reported would infer that
the remaining defaults --- \texttt{norm='l2'} among them --- were in force in the library
version the article names \citep{pedregosa2011}. We make that inference and we hold it to be
reasonable. Having relied on it here, we do not also count the silence on \texttt{norm} as a
reproducibility failure: \S\ref{sec:original-consequence} rests that charge on the
aggregation rule, which is a modelling step rather than a library setting and which no
default convention can supply, on the unpublished keyword list, and on the contradiction
between the original's prose and its printed formula on score normalisation. Those three are
unanswerable from any convention; this one is not, and we say so in both places.

The inference is nonetheless an \emph{inference}, it is labelled as one wherever it is used,
and it can be denied at no cost to the argument: what follows is a property of $\suslag$ ---
our reconstruction under that inference (\S\ref{sec:index}) --- and of any pipeline built the
same way. It is not a documented choice of the original, and we do not assert that the
original made it.

\subsection{Vector normalisation discards magnitude}

Vector normalisation --- L2 scaling of a document's term vector to unit length, applied
within the document at the TF-IDF stage --- retains the vector's direction, that is the
\emph{mix} across terms, and discards its length, that is \emph{how much} there is. The
algebra is immediate: for any $\mathbf{v} \neq \mathbf{0}$ and any scalar $c > 0$,

\begin{equation}
\frac{c\,\mathbf{v}}{\lVert c\,\mathbf{v}\rVert_2}
= \frac{c\,\mathbf{v}}{c\,\lVert \mathbf{v}\rVert_2}
= \frac{\mathbf{v}}{\lVert \mathbf{v}\rVert_2}.
\label{eq:l2-invariance}
\end{equation}

Two documents whose ESG term counts are proportional to one another therefore receive
\emph{identical} unit vectors. Nothing downstream recovers the difference: the aggregation
from per-term values to a document score is then applied to identical inputs and returns
identical outputs, whatever that aggregation is --- sum, mean, weighted mean or anything
else. This is not an approximation or a small-sample effect; it is an identity.

The distinction matters because it separates the two silences in the original's
specification, which are unequal in both respects that matter. They are unequal in what a
convention can close: \texttt{norm} has a library default and the aggregation rule has none
(\S\ref{sec:original}). They are also unequal in consequence, and in the opposite direction.
Sum and mean differ by the constant factor $1/V$, and $\Z(\cdot)$ annihilates a positive
affine rescaling of a component, so on our index the aggregation choice cannot move a single
value (\S\ref{sec:index}) --- the silence no convention closes turns out to be the harmless
one here. The vector normalisation is not of that kind. It does not rescale the score; it
changes what the score is a function of, which is why an inference, however reasonable, is a
thin thing to leave a first-order choice resting on.

Whether that is a defect depends on the construct being measured. Discarding magnitude is
appropriate when the question is which document best matches a query. It is not appropriate
here: $\ESGSI$ contrasts the tone of a report against the amount of substance behind it, so
a substance component that returns the same value for a document and for the same document
carrying twenty times the ESG content is measuring on the wrong side of that comparison.

\subsection{Two demonstrations}

Equation~\eqref{eq:l2-invariance} predicts two failures. Both are executed end to end on the
live pipeline rather than argued, so that the pipeline is shown to behave as the algebra
says rather than assumed to.

\emph{Amplification.} We take a document synthesised from the observed vocabulary counts of
one report --- synthesised so that the counts can be scaled exactly without perturbing the
$n$-gram structure --- multiply those counts by $k \in \{1, 2, 5, 20\}$, re-embed the result
in the remaining reports and re-run the full vectorisation. Scaling a document's own counts
leaves every $df_j$ unchanged, since a term present remains present and a term absent remains
absent, so the idf vector is fixed across the four runs and the comparison is clean.
Table~\ref{tab:measurement-scaling} reports the outcome. Twenty times the ESG content leaves
$\suslag$ unmoved to machine precision, while any specification with a length denominator
rises by the factor $k L_1 / L_k$, as it should.

\begin{table}[ht]
\centering
\caption{The same document with ESG term counts multiplied by $k$. $\suslag$ is reported
relative to its value at $k = 1$: Equation~\eqref{eq:l2-invariance} does not depend on the idf
vector, so the ratio is exact for any corpus, whereas the level is corpus-dependent. The
maximum absolute deviation across the four runs is $6.9 \times 10^{-18}$.}
\label{tab:measurement-scaling}
\begin{tabular}{rrrrr}
\toprule
$k$ & Tokens & ESG hits & $\susdens$ & $\suslag$, relative to $k=1$ \\
\midrule
1  & 16{,}556 & 700      & 4.2281  & 1.000000 \\
2  & 18{,}112 & 1{,}400  & 7.7297  & 1.000000 \\
5  & 22{,}780 & 3{,}500  & 15.3644 & 1.000000 \\
20 & 46{,}120 & 14{,}000 & 30.3556 & 1.000000 \\
\bottomrule
\end{tabular}
\end{table}

\emph{Dilution.} The converse is the case a reader will recognise from practice. We append
56{,}000 non-ESG filler tokens, drawn from the document's own frequent non-vocabulary words,
to a 22{,}576-token report, leaving its 695 ESG mentions untouched.
Table~\ref{tab:measurement-dilution} reports the outcome: measured ESG density falls from
3.0785 to 0.8845 per 100 tokens, a $3.48$-fold dilution, and $\suslag$ does not move to ten
decimal places. A report can be padded without limit --- with boilerplate, with repetition,
with any material outside the ESG vocabulary --- at no cost to its measured sustainability
substance.

\begin{table}[ht]
\centering
\caption{Dilution of a report with non-ESG filler text.}
\label{tab:measurement-dilution}
\begin{tabular}{lrrrr}
\toprule
Variant & Tokens & ESG hits & $\susdens$ & $\suslag$ \\
\midrule
Original             & 22{,}576 & 695 & 3.0785 & 0.0347763839 \\
$+\,56{,}000$ filler & 78{,}576 & 695 & 0.8845 & 0.0347763839 \\
\bottomrule
\end{tabular}
\end{table}

\subsection{What survives vector normalisation is a spread measure}

If magnitude is discarded, something else is being measured, and it is worth naming. Let
$\mathbf{u} \geq \mathbf{0}$ be a unit vector in $\mathbb{R}^{V}$, with $V = 154$ the
vocabulary size (\S\ref{sec:index}). Then

\begin{equation}
1 \;\leq\; \sum_{j=1}^{V} u_j \;\leq\; \sqrt{V}.
\label{eq:unit-sum-bounds}
\end{equation}

In Equation~\eqref{eq:unit-sum-bounds} the lower bound follows from
$\bigl(\sum_j u_j\bigr)^2 = \sum_j u_j^2 + \sum_{j \neq l} u_j
u_l \geq \lVert \mathbf{u} \rVert_2^2 = 1$ for non-negative $\mathbf{u}$, with equality when
all mass sits on a single term; the upper bound is Cauchy--Schwarz, with equality when the
mass is spread evenly over all $V$ terms. Between the two, mass spread evenly over $m$ terms
gives $\sum_j u_j = \sqrt{m}$. Aggregating an L2-normalised vector over the vocabulary ---
$\suslag$ is that sum divided by $V$ --- therefore rewards \emph{dispersion across terms},
not \emph{amount of content}: a report mentioning many distinct ESG topics once each scores
above a report treating a few of them intensively, and the ceiling is set by the size of the
vocabulary rather than by the disclosure.

That is what the data show. On our corpus $\suslag$ correlates $0.9733$ with $\BREADTH$, the
effective number of distinct vocabulary terms (\S\ref{sec:index}), and $0.4177$ with ESG
density. It is, to a correlation of $0.97$, a measure of topical breadth carrying the name of
a measure of disclosure intensity. Breadth is a legitimate construct and may well be worth
reporting; it is simply not the construct the index is defined to capture.

\subsection{The term weighting is second-order}

Inverse document frequency down-weights terms that appear everywhere, because in retrieval a
term present in every document separates nothing. In a corpus of ESG disclosure the same
terms are present everywhere for the opposite reason: they are what the genre is about.
Rarity and topical centrality are inversely related here, so idf transfers weight away from
the vocabulary that carries the disclosure and towards the vocabulary that does not.
Table~\ref{tab:measurement-idf} gives the inversion on our corpus.

\begin{table}[ht]
\centering
\caption{Concentration of ESG term occurrences by idf band ($n = 343$).}
\label{tab:measurement-idf}
\begin{tabular}{lrr}
\toprule
Band & Terms & Share of all ESG occurrences \\
\midrule
$\mathrm{idf} < 1.2$ (document frequency $> 280$) & 43 & 78.35\,\% \\
$\mathrm{idf} > 2$                                & 46 & 2.80\,\% \\
\bottomrule
\end{tabular}
\end{table}

Forty-three terms hold 78.35\,\% of all ESG occurrences in the corpus and are weighted at
close to the floor of $1$; forty-six hold 2.80\,\% and are weighted at up to nearly five
times it. The exchange rate is explicit at the level of individual terms. \emph{Fair
condition} occurs eight times in the entire corpus and carries $\mathrm{idf} = 4.8947$;
\emph{climate}, among the most frequent entries in the vocabulary, carries
$\mathrm{idf} = 1.0088$. One mention of the former is therefore worth $4.852$ mentions of the
latter. No account of what a sustainability report discloses supports that exchange rate, and
no reader of such a report would recognise it.

It is nonetheless the weaker of the two criticisms, and we say so rather than present the
pair as equals. The inversion is theoretically indefensible and empirically close to inert,
for the reason the table itself supplies: the terms idf promotes carry 2.80\,\% of the mass.
Holding the treatment of length fixed --- $\susdens$ and $\sustfidf$ both divide by document
length, and neither applies vector normalisation (\S\ref{sec:index}) --- and varying only the
term weighting, the two specifications correlate $0.9926$, and their assembled indices
disagree about the label of $6$ of $343$ reports. Varying the vector normalisation instead
moves $96$ of $343$ (\S\ref{sec:robustness}). A reader who accepts the first-order argument and
rejects this subsection in its entirety loses nothing that matters here.

The concession runs further, and stating it is what keeps the first-order claim credible.
Nothing above indicts idf weighting as such for measurement use: with the vector
normalisation replaced by a length denominator, the same idf weights produce a specification
that agrees with our independent yardstick at least as well as any other we compute
(\S\ref{sec:validity}). What this section establishes is a defect of one library default, not
of the weighting scheme and not of the family.

\subsection{The sentiment component is exposed to the same class of problem}

The critique above concerns one half of the index. The other half is $\SEN_i$, TextBlob
polarity, which enters $\ESGSI$ with the same weight as the substance component
(\S\ref{sec:index}). TextBlob's polarity is a pattern-matching lexicon whose scores are
modified by negation and by intensifiers. The original's preprocessing, disclosed in its
Section~3.1.1 (p.~3), lowercases, strips URLs, numerals and punctuation, removes NLTK's
English stopwords augmented by an unpublished custom list, discards every token shorter than
three characters, and lemmatises \citep{lagasio2024}. That pipeline removes the tokens those
rules operate on --- \emph{no} falls below the three-character floor, \emph{not} and
\emph{very} are NLTK stopwords --- and stopword removal breaks the adjacency the modifier
rules read. A polarity score is still returned; it is no longer the quantity the instrument
was validated to produce. Whether the raw text or this preprocessed output reached TextBlob
is never stated (\S\ref{sec:original}), and the five-fold cross-validation announced for the
sentiment model is never reported. The instrument is therefore either applied outside its
validity conditions or it is not, and the article does not say which.
Our own substitution, a count-based finance dictionary that presupposes no negation or
intensifier parsing (\S\ref{sec:index}), is a divergence by design and not a repair.

\subsection{What carries into the next section}

Both criticisms share a structure. There is a vector-normalisation choice at the TF-IDF stage
and a score-normalisation choice at index assembly; at each stage that choice dominates the
weighting choice; and the original states neither (\S\ref{sec:original}). The two are not
undisclosed in the same way, and the difference cuts against us in the first case and for us
in the second: the first is recoverable from a reporting convention, as conceded above,
whereas the second is not merely unstated but stated twice and inconsistently, in prose and
in a printed formula that specify different estimators. The second stage is the subject of
\S\ref{sec:artefact}. What the present section supplies is the reason to expect a result
there: once first-order choices of this kind are left to a convention or to a contradiction,
a published conclusion that turns on one of them is a predictable outcome rather than a
coincidence.
